\section{Related Work}
\label{sec:related}

\paragraph{Fabric and surface defect detection.}
Automated fabric inspection has moved from hand-crafted texture analysis to deep
detectors~\cite{fabricsurvey}. Recent fabric-specific models improve sensitivity
to small and subtle defects through fine-grained convolutions, attention
modules~\cite{cbam}, small-object detection layers, and tiny-object loss
functions~\cite{fabricyolo,liang2025denim}; recent surveys review this rapidly
growing literature~\cite{eaaisurvey}. These methods frame the residual
difficulty as a model-capacity problem and seek a stronger low-resolution
detector; they do not separate misses that are recoverable in principle from
 misses caused by insufficient input resolution.

\paragraph{Small-object and industrial defect detection.}
Small-object detection is a long-standing challenge, addressed through
higher-resolution features, scale-aware training, context, and
augmentation~\cite{sodsurvey,focalloss}. In industrial inspection, benchmarks
span metallic surface defects~\cite{neudet} and unsupervised anomaly
detection~\cite{mvtecad}, and class-imbalance-aware detectors have been
proposed for defects with rare categories~\cite{eaaiimbalance}. Our work is complementary: rather than proposing a
stronger detector or anomaly model, we analyze when input resolution is the
binding constraint and act on that diagnosis.

\paragraph{Object detectors.}
Modern object detection spans two-stage detectors built on region
proposals~\cite{fasterrcnn} and single-stage detectors that regress boxes
directly~\cite{yolov1,ssd}, typically over deep residual
backbones~\cite{resnet} with multi-scale feature
pyramids~\cite{fpn}. We build on standard real-time detectors spanning the
anchor-free YOLO line~\cite{yolov1,yolov8,yolo11} and the
detection-transformer line (DETR~\cite{detr}, Deformable DETR~\cite{deformabledetr}, RT-DETR~\cite{rtdetr}); our
diagnosis and framework treat the detector as a frozen black box and are not
specific to either family.

\paragraph{Resolution, tiling, and adaptive inference.}
High input resolution is a well-known lever for small-object detection, and
sliding-window or tiling schemes such as SAHI~\cite{sahi} process an image at high
resolution in pieces. These apply high resolution \emph{uniformly} (to the whole
image or to a regular grid), incurring large cost. Adaptive-resolution and
coarse-to-fine methods~\cite{ranet,drnet,dynamiczoom} reduce this cost by
selecting where to look or at what resolution to process, typically driven by
objectness, confidence, or a learned resolution predictor. VBAD differs in its supervision:
the trigger is learned from \emph{cross-resolution failure modes}---images where
low resolution demonstrably fails and high resolution demonstrably recovers---
rather than from generic saliency or confidence heuristics, tying the decision
directly to when high resolution is actually worthwhile.

\paragraph{Class imbalance and long-tailed recognition.}
Fabric defect categories are heavily imbalanced. Long-tailed recognition has
been tackled by loss reweighting~\cite{eql,focalloss}, decoupling
representation from classifier learning~\cite{decoupling}, bilateral-branch
architectures~\cite{bbn}, and minority-class augmentation such as
copy-paste~\cite{copypaste}. Our diagnosis shows, however, that on this fabric
benchmark per-class difficulty is governed by visibility rather than frequency,
so frequency-indexed remedies are not the binding lever; we therefore focus on
the resolution dimension.

\paragraph{Detector error analysis.}
Diagnostic toolkits such as TIDE~\cite{errortaxonomy} decompose detector errors
into interpretable types. We instantiate a resolution-centric analysis for
industrial inspection and, unlike purely diagnostic work, turn the analysis into
an operational triage policy and a reliability boundary.
